\documentclass[letterpaper,12pt]{article}
\usepackage[scale=0.9]{geometry}

\usepackage{url}
\usepackage{parskip}
\RequirePackage{color}
\usepackage[usenames,dvipsnames]{xcolor}
\definecolor{linkcolour}{rgb}{0,0.2,0.6}

\usepackage{tabularx}
\newcolumntype{C}{>{\centering\arraybackslash}X}

\usepackage{enumitem}
\setlist[itemize]{label=$\cdot$,nosep}

\usepackage{titlesec}
\titleformat{\section}{\Large}{}{0em}{}[\titlerule]
\titlespacing{\section}{0pt}{6pt}{6pt}

\usepackage[style=authoryear,sorting=ddnt, maxbibnames=2, backend=biber]{biblatex}

\usepackage[unicode, draft=false]{hyperref}
\hypersetup{colorlinks,breaklinks,urlcolor=linkcolour,linkcolor=linkcolour}
\addbibresource{citations.bib}

\usepackage{fontawesome5}

\newenvironment{job}[2]{
	\textbf{#1} \hfill #2\par\vspace{-0.5\baselineskip}
}{}

\begin{document}
\pagestyle{empty}

\begin{center}
	\Huge{Youqiu Zhan}\\
	\normalsize
	\href{https://ulysseszh.github.io}{\raisebox{-0.05\height}\faGlobe~ulysseszh.github.io}\qquad
	\href{mailto:Ulysses Zhan <UlyssesZhan@gmail.com>}{\raisebox{-0.05\height}\faEnvelope~UlyssesZhan@gmail.com}
\end{center}

\section{Education and coursework}

\begin{job}{University of California, Santa Barbara}{2022--2026 (expected)}
	Bachelor of Science in physics\hfill(GPA: 3.92/4.0)\par\vspace{-0.5\baselineskip}
	\begin{itemize}
		\item Quantum many-body problem (PHYS 217A--B)\hfill A, A+\\
		Final projects:
		\href{https://ulysseszh.github.io/coursework/crust-neutron-star.pdf}{Quantum many-body theories about the crust of neutron stars},
		\href{https://ulysseszh.github.io/coursework/axion-insulator.pdf}{Axion electrodynamics in topological insulators}
		\item Relativistic quantum field (PHYS 221B--C)\hfill A, audit
		\item Condensed matter (PHYS 223A--C)\hfill A+, A\textminus, A\\
		Final project:
		\href{https://ulysseszh.github.io/coursework/exact-2d-percolation.pdf}{Exact results in 2D percolation}
		\item Gauge theory (PHYS 229A--B)\hfill audit
		\item General relativity (PHYS 231A--C)\hfill audit, A, A\\
		Final project:
		\href{https://ulysseszh.github.io/coursework/spinor-curved.pdf}{Spinor fields in curved spacetime}
	\end{itemize}
\end{job}

\begin{job}{High School Affiliated to Shanghai Jiao Tong University}{2019--2022}
	International Baccalaureate\hfill(Grade: 42/45)
\end{job}

\section{Research experience}

\begin{job}{Bouwmeester Lab}{2024--}
\begin{itemize}
	\item Using COMSOL to simulate phononic crystals on a thin plate to find bind gaps
	and design ultra-high-$Q$ mechanical resonators.
	\item Using COMSOL to simulate third sound modes of superfluid helium on 2D geometries.
	\item Using KLayout to design and produce masks used in nanofabrication of mechanical devices
	used in optomechanical experiments.
	\item Using QuTiP to simulate and visualize optomechanical processes
	such as stimulated adiabatic Raman passage.
	\item Using QuTiP to simulate quantum models of third sound modes of superfluid helium on 2D geometries.
\end{itemize}
\end{job}

\begin{job}{UCSB honor thesis}{2025--}
	Studying the continuum limit of statistical problem of the distribution of the maximum run length in
	a Bernoulli sequence using techniques from mathematical physics,
	such as the Adomian decomposition method and the grand canonical ensemble.
	Using Monte Carlo methods to simulate the discrete problem to verify theoretical results.
\end{job}

\begin{job}{Independent research}{2021}
	Using PyTorch to implement and train self-supervised neural networks
	for finding the Hamiltonian of classical mechanical systems
	from trajectory data.
\end{job}

\begin{job}{Regeneron ISEF}{2020--2021}
	Using web technology to simulate and visualize classical Hamiltonian mechanical systems.
	Implementing integrators and signal processing algorithms in JavaScript.
\end{job}

\begin{job}{Independent research}{2020--2021}
	Using ABAQUS to simulate and visualize stress distribution in 3D structures
	fatigue life under cyclic loading conditions.
\end{job}

\section{Development projects}

Individual projects:\par\vspace{-0.5\baselineskip}
\begin{tabularx}{\linewidth}{lX}
\href{https://ulysseszh.github.io/doc/alda-rb/}{alda-rb} &
Ruby library and DSL for live-coding music, powered by Alda \\
\href{https://github.com/UlyssesZh/DisplayToggle}{DisplayToggle} &
Android app for controlling the power mode of the physical display \\
\href{https://github.com/UlyssesZh/mechsimul-js}{mechsimul-js} &
Hamiltonian mechanics simulator written in JavaScript \\
\href{https://sunniesnow.github.io/game/}{Sunniesnow} &
Web rhythm game inspired by Lyrica \\
\href{https://github.com/UlyssesZh/SunProxy}{SunProxy} &
Android app for setting up network packet redirection rules by VPN \\
\end{tabularx}

Contributed projects:\par\vspace{-0.5\baselineskip}
\begin{tabularx}{\linewidth}{lX}
\href{https://phil294.github.io/AHK_X11}{AHK\_X11} &
AutoHotkey for X11-based systems \\
\href{https://github.com/nixos/nixpkgs}{Nixpkgs} &
Software package collection distributed with Nix \\
\href{https://heerdebeer.org/Software/markdown/paru/}{Paru} &
Ruby wrapper over Pandoc \\
\href{https://pixijs.com}{PixiJS} &
2D rendering engine built on WebGL and WebGPU \\
\end{tabularx}

\section{Publications}
\begin{refsection}[citations.bib]
\nocite{*}
\printbibliography[heading=none]
\end{refsection}

\section{Skills}
\begin{itemize}
	\item Programming languages: Ruby, JavaScript (and TypeScript),
	Java (and Kotlin), Python
	\item DevOps and software distribution: GNU/Linux, Docker, SSH, Git, Nix
	\item Development tech stacks: Android SDK, Node.js, Electron.js, Jekyll, Pandoc
	\item Scientific software: COMSOL, KLayout, BEAMER, Wolfram, Praat
	\item Scientific Python packages: QuTiP, NetKet, PyTorch, OpenCV
	\item Spoken languages: English, Mandarin Chinese
\end{itemize}

\vfill
\center{\footnotesize Last updated: \today}

\end{document}
